\documentclass[../main.tex]{subfiles}
\begin{document}

% ---------------------------------------------------------------------------
% Capa institucional UFMG
% ---------------------------------------------------------------------------
\begin{frame}[plain]
  \begin{tikzpicture}[overlay,remember picture]
    \node[left=-0.15cm] at (current page.0){
      \includegraphics[scale=0.202]{\CaminhoCapa}
    };
  \end{tikzpicture}
  \titlepage
\end{frame}

% ---------------------------------------------------------------------------
% Agenda
% ---------------------------------------------------------------------------
\section[Agenda]{}

\begin{frame}
  \frametitle{Agenda}
  \tableofcontents
\end{frame}

% ---------------------------------------------------------------------------
% Motivacao e Objetivos
% ---------------------------------------------------------------------------
\begin{frame}
  \frametitle{Motiva\c{c}\~ao e Objetivos na Cora Z7-10}
  \begin{columns}[T]
    \begin{column}{.52\textwidth}
      \begin{itemize}
        \item \textbf{Contexto de borda:} enlaces restritos exigem compress\~ao e cifragem antes da transmiss\~ao \autocite{NISTSP80038A}.
        \item \textbf{Gargalo em CPU:} executar as duas rotinas apenas em software consome ciclos do ARM Cortex-A9.
        \item \textbf{Plataforma alvo:} SoC Zynq-7010 (\texttt{XC7Z010}) na Digilent Cora Z7-10 (PS a 650 MHz e PL a 100 MHz) \autocite{DigilentCoraZ7,XilinxUG585}.
        \item \textbf{Est\'imulo f\'isico:} microcontrolador ESP32 injeta quadros sint\'eticos via SPI (1035 B) para valida\c{c}\~ao HIL.
      \end{itemize}
    \end{column}
    \begin{column}{.44\textwidth}
      \begin{pontoschave}[Objetivos do Projeto]
        \begin{itemize}
          \item Avaliar o pipeline \textit{Compress-then-Encrypt} (\textbf{LZSS + AES-128-CTR}).
          \item Comparar 4 cen\'arios de particionamento PS/PL com carga fixa de \textbf{1024 B} (ADR-020).
          \item Medir lat\^encia real via ARM PMU (\texttt{PMCCNTR}, 28 ciclos de overhead) e contadores na PL.
          \item Caracterizar ocupa\c{c}\~ao de sil\'icio (LUTs, FFs, BRAMs) e fechamento temporal.
        \end{itemize}
      \end{pontoschave}
    \end{column}
  \end{columns}
  \fonte{README.md, STATUS.md, ADR-020}
\end{frame}

% ---------------------------------------------------------------------------
% Estrutura da Apresentacao
% ---------------------------------------------------------------------------
\begin{frame}
  \frametitle{Estrutura da Apresenta\c{c}\~ao}
  \begin{roteiro}
    \linharoteiro{Se\c{c}\~ao 1}{Motiva\c{c}\~ao, objetivos e plataforma Zynq-7010}{Cora Z7-10 + ESP32}
    \linharoteiro{Se\c{c}\~ao 2}{Topologia PS/PL, piv\^o ADR-012/017 e BRAM Dual-Port}{Interconex\~ao AXI/BRAM}
    \linharoteiro{Se\c{c}\~ao 3}{Coprocessadores VHDL-2008, firmware C11 e injetor SPI}{RTL, PS e ESP32}
    \linharoteiro{Se\c{c}\~ao 4}{Matriz ADR-020, resultados HIL, FPGA e Risco R9}{M\'etricas Emp\'iricas}
    \linharoteiro{Se\c{c}\~ao 5}{S\'intese PS vs. PL vs. H\'ibrido e trabalhos futuros}{Conclus\~ao}
  \end{roteiro}
  \par\vspace{1.5ex}
  \begin{alertblock}{Quest\~ao Central de Co-Projeto}
    \small Em blocos de telemetria de 1024 B, quais etapas do pipeline compensam o custo de transfer\^encia PS/PL via BRAM?
  \end{alertblock}
  \fonte{README.md, Matriz\_de\_Testes\_e\_M\'etricas.csv}
\end{frame}

\end{document}
